\section{Síntesis y ampliación}

Esta sección condensa todo el episodio en una hoja de ruta. La evolución de la conducción autónoma a lo largo de diez años puede verse como tres «aumentos de dimensión»: el primero va de las reglas y los mapas de alta precisión a la representación espacial unificada con BEV/Transformer; el segundo, del desarrollo modular de funciones al entrenamiento de capacidades de extremo a extremo; el tercero, de la imitación de las trayectorias de conducción humanas al Agent del mundo físico basado en VLM, modelos del mundo y RL.

\begin{importantbox}{Seis conclusiones centrales}
Primera: la ruta de mapas de alta precisión y LiDAR resuelve la certidumbre a corto plazo, pero difícilmente escala al nivel de ciudades enteras. Segunda: la Camera tiene alta densidad de información y bajo costo, y es la base de la ruta de visión para producción en serie. Tercera: lo esencial de BEV es unificar la representación espacial de varias cámaras y reducir los errores de fusión tardía. Cuarta: el enfoque de extremo a extremo representa un cambio de paradigma, del desarrollo de funciones al entrenamiento de capacidades. Quinta: la esencia de L3/L4 es el límite de responsabilidad, no una lista de funciones. Sexta: la barrera de largo plazo de la conducción autónoma en producción en serie proviene de los datos reales, la capacidad de cómputo, la ejecución organizacional y el ciclo cerrado de validación de seguridad.
\end{importantbox}

\subsection{Tabla de consulta rápida de términos clave}

\noindent\begin{tabularx}{\textwidth}{p{0.22\textwidth}p{0.34\textwidth}X}
\toprule
Término & Significado en este episodio & Relación con el curso \\
\midrule
Mapa de alta precisión & Registro previo de la geometría y la información semántica de las vías & Dependencia central de la primera ruta basada en certidumbre. \\
LiDAR & Sensor activo de medición de distancia & Da geometría directa, pero tiene restricciones de costo y de densidad de información. \\
Camera & Sensor de visión pasivo & Alta densidad de información; adecuado para producción en serie de bajo costo y datos a gran escala. \\
BEV & Bird's Eye View, representación espacial a vista de pájaro & Paso clave para el modelado unificado de varias cámaras. \\
Transformer & Arquitectura general de modelado de secuencias/atención & Se trasladó a la comprensión espacial en la conducción autónoma. \\
De extremo a extremo & Mapeo modelado desde la entrada hasta la trayectoria/acción & La conducción autónoma pasa de la ingeniería de funciones al entrenamiento de capacidades. \\
VLM & Vision-Language Model, modelo de visión y lenguaje & Introduce comprensión semántica e interpretación de escenarios complejos. \\
World Model & Modelo que predice las consecuencias de las acciones y el estado futuro del mundo & Sustenta la planificación, la simulación y el ciclo cerrado de RL. \\
RL & Reinforcement Learning, aprendizaje por refuerzo & Optimiza la política con retroalimentación; es una de las direcciones de la conducción autónoma en ciclo cerrado. \\
\bottomrule
\end{tabularx}

\subsection{Preguntas para seguir observando}

\begin{enumerate}
\item ¿Podrá la ruta de extremo a extremo superar de forma estable a los sistemas basados en reglas en más ciudades, condiciones climáticas, estructuras viales y escenarios de cola larga?
\item ¿Cómo se reequilibrarán las rutas Camera-only, LiDAR-heavy y de fusión multisensor entre costo, seguridad y regulación?
\item ¿Se convertirá el modelo del mundo en el núcleo de la siguiente etapa de la conducción autónoma, o solo servirá como apoyo para entrenamiento/simulación?
\item Cuando L3 llegue a la producción en serie en China, ¿cómo deben diseñarse el límite de responsabilidad, el mecanismo de toma de control y la educación del usuario?
\item ¿Puede el ciclo cerrado de datos de flotilla de los fabricantes de automóviles convertirse en una barrera de largo plazo difícil de replicar para las empresas de modelos?
\item ¿Cómo equilibrar la alta presión organizacional y la validación de seguridad para no sacrificar la confiabilidad en aras de la velocidad de lanzamiento?
\end{enumerate}

\subsection{Lecturas adicionales}

\begin{itemize}
\item Si le interesan los modelos de negocio de la conducción autónoma y las diferencias entre Waymo y Momenta, puede contrastarlo con la entrevista del EP132 a Gao Jiyang (高继扬), de Xinghaitu (星海图).
\item Si le interesan la Physical AI y la fábrica de modelos grandes a bordo del vehículo, puede contrastarlo con la entrevista del EP120 a Liu Xianming (刘先明), de Xiaopeng (小鹏).
\item Si le interesan VLA, los modelos del mundo y los modelos fundacionales para robótica, puede contrastarlo con la versión de pantalla compartida sobre VLA `eiQFomOuCJs` y con la entrevista sobre multimodalidad del EP102 a Zhang Xiangyu (张祥雨).
\item Si le interesan los costos de hardware y la industria del LiDAR, puede contrastarlo con la entrevista del EP111 a Hesai (禾赛).
\end{itemize}

\end{document}
